\section{Introduction}\label{sec:introduction}

A real Littlewood polynomial of length $N$ is a polynomial
\[
 P(z)=\sum_{k=0}^{N-1}\varepsilon_k z^k,
 \qquad \varepsilon_k\in\{-1,1\}.
\]
Its squared mean modulus on the unit circle is $N$, by Parseval's
identity. Thus $\sqrt N$ is the natural scale, and its maximum modulus
is at least $\sqrt N$. A family is \emph{ultraflat} if
$|P(z)|/\sqrt N$ tends uniformly to $1$ on the unit circle as the
length tends to infinity. We prove that real Littlewood polynomials
can be ultraflat, through every sufficiently large integer length.

\begin{theorem}\label{thm:main}
For every $\varepsilon\in(0,1)$ there is an integer $N_0$ such that,
for every integer $N\ge N_0$, there are signs
$\varepsilon_0,\ldots,\varepsilon_{N-1}$ for which
\[
 (1-\varepsilon)\sqrt N
 \le\left|\sum_{k=0}^{N-1}\varepsilon_k z^k\right|
 \le(1+\varepsilon)\sqrt N
 \qquad\text{for every }|z|=1.
\]
\end{theorem}

The signs may depend on $N$. The theorem concerns both extrema of the
same polynomial on the entire circle, including the real points
$z=1$ and $z=-1$.

\subsection{The flatness problem}

Erd\H{o}s asked whether real-sign polynomials could stay between two
fixed positive multiples of their square-root scale on the entire
circle \cite[Problem~26]{Erdos1957}; the problem also appears in
Littlewood's work \cite{Littlewood1966}. The Rudin--Shapiro construction
gives real-sign polynomials with maximum at most $\sqrt{2N}$ at dyadic
lengths \cite{Rudin1959}. Balister, Bollob\'as, Morris, Sahasrabudhe and
Tiba proved the two-sided constant-factor statement at every degree
at least two \cite[Theorem~1.1]{BalisterEtAl2020}. Their theorem controls
the scale of both extrema, while ultraflatness requires both normalized
extrema to converge to one.

For complex coefficients of modulus one, Kahane constructed ultraflat
polynomials \cite{Kahane1980}. Bombieri and Bourgain improved the
quantitative bounds and gave effective constructions
\cite[Theorems~4 and~7]{BombieriBourgain2009}. The restriction to real
signs is substantial. The predecessor \cite[Theorem~1.1]{OpenAIMinimal}
obtains an asymptotically minimal maximum through all lengths. Its
version-2 refinement \cite[Theorem~1.1]{OpenAILowerEnvelope2026} gives
the simultaneous bounds
\[
 \frac1{16}\sqrt N\le |P(z)|\le(1+\eta)\sqrt N
 \qquad(|z|=1)
\]
for every fixed $\eta>0$ and all sufficiently large $N$.
Theorem~\ref{thm:main} strengthens that lower bound to its asymptotically
optimal value. The present paper develops the additional constructions
needed for this strengthening and uses the retained version-2 paper
as the source of two shared lemmas.

Recent obstructions distinguish the relevant quantifiers. Erd\'elyi
proved that the full sequence of partial sums of one fixed unimodular
power series cannot be ultraflat
\cite[Theorem~2.1]{ErdelyiSections2026}. This does not restrict the
independent choice of signs at each length in Theorem~\ref{thm:main}.
His separate lower bound
\cite[Theorem~2.1]{Erdelyi2026},
$\max_{|z|=1}|P(z)|^2\ge N+(N-1)^{1/3}/38$, is compatible with
the asymptotic conclusion here.\footnote{Contrary nonflatness claims
in \cite{ElAbdalaoui2025,ElAbdalaoui2025Generalization} and specific
issues in their arguments are discussed in
\cite[Appendix~A]{OpenAIMinimal}. Those claims are not inputs to the
present proof.}

\subsection{Construction and relation to the earlier version}

The proof first constructs a continuous function $B_N$ on the circle
whose modulus is close to one. Its Fourier coefficients of indices
$0,\ldots,N-1$ are real and at most $(1+o(1))/\sqrt N$ in absolute
value, and the sum of the absolute values of all remaining coefficients
tends to zero. Parseval's identity then forces the retained coefficients
to be almost of sign size in aggregate. A discrepancy estimate rounds
them to signs with a small uniform error. This last passage converts
the continuous construction into a Littlewood polynomial.

To build $B_N$, we adapt the auxiliary-torus and phase-correction
arguments of \cite[Sections~2, 4, and~5]{OpenAILowerEnvelope2026}.
The auxiliary object is a real trigonometric polynomial
$F(y)=\sum_{a\in\mathcal A}c_a\e(a\cdot y)$ on a higher-dimensional
torus, where $\e(t)=\exp(2\pi i t)$. Its norm is at most $1+\delta$,
and its Fourier coefficients
are compared with positive weights $w_a=|a\cdot v|$ whose sum is
close to one. Here $\delta>0$ is a small error parameter and $v$ is
a fixed real vector. The first additional requirement is the two-sided
comparison
\[
 1\le |c_a|/\sqrt{w_a}\le1+C\delta.
\]
We obtain it by spreading each original coefficient over a box of
frequencies and rounding the spreading factors to equal coefficient
magnitudes. Section~\ref{sec:rounding} proves the required complex
version of the defect-sensitive rounding estimate, and
Section~\ref{sec:balanced} carries out the auxiliary construction.

The frequencies of $F$ then supply oscillatory waves on disjoint short
arcs, each with constant modulus $|c_a|/\sqrt{w_a}$. Their phases are
chosen so that, at a fixed Fourier index, the leading stationary-phase
terms sum to a single value of $F$ multiplied by a factor in $[0,1]$.
The norm bound for $F$ therefore controls these contributions without accumulating
the magnitudes of its coefficients.

The second change concerns the endpoints of the waves. We increase the
magnitude of their phase curvature near the endpoints while keeping
their amplitudes constant. Stationary phase
therefore gives small endpoint Fourier contributions while the
function's modulus stays close to one. The waves are joined across
the remaining gaps by matching their values and their leading phase
derivatives. These matching conditions cancel the first boundary terms
in integration by parts and give a summable exterior Fourier tail.
Section~\ref{sec:oscillation} supplies the analytic estimates,
Section~\ref{sec:waves} constructs and joins the waves, and
Section~\ref{sec:completion} projects and rounds to signs.

The signed interval packing and real matrix discrepancy lemmas are
inherited from \cite[Lemmas~3.1 and~6.2]{OpenAILowerEnvelope2026}.
Their combinatorial inputs are the Pippenger--Spencer hypergraph-coloring
theorem \cite{PippengerSpencer1989}, in the quantified form of
\cite[Lemma~2.1]{AlonYuster2005}, and the partial-coloring method of
Spencer \cite{Spencer1985} and Lovett--Meka \cite{LovettMeka2015}.
Discrepancy and the repair of exceptional arcs also occur in
\cite[Sections~2.1 and~4]{BalisterEtAl2020}, where a sine polynomial
corrects a Rudin--Shapiro cosine polynomial. Here the correction joins
oscillatory waves while preserving their almost constant modulus.
All dimensions, packing data, intervals, and leading phases are fixed before
the final length tends to infinity; the construction consequently
applies to all sufficiently large integer lengths.
